\scanpage{52}
\mspage{\handwritten{20}\,\raisebox{.5ex}{\begin{tikzpicture}[x=1cm,y=1cm,baseline=0pt]\draw[line width=.6pt] (0,.04)--(.025,0)--(0,-.025);\end{tikzpicture}}}
% archive sheet number: 51

\deleted{Let $\handwritten{\varphi}_i:G\handwritten{\longrightarrow}\text{\uline{GL}}_k(V_i)$ denote the morphism determined by the action of $G$ on $V_i$, and $\handwritten{\varphi}:G\handwritten{\longrightarrow}\text{\uline{GL}}_k(A)$ the morphism determined by the action of $G$ on $A$. Then, since $A$ is the union of the $V_i$, $\operatorname{Ker}\handwritten{\varphi}$ is the \unsure{intersection|intersection} of the $\operatorname{Ker}\handwritten{\varphi}_i$; now $\operatorname{Ker}\handwritten{\varphi}$ is the trivial $k$-group; since $G$ is n\oe{}therian, there exists $i\in I$ such that $\operatorname{Ker}\handwritten{\varphi}_i$ is the trivial $k$-group, which implies (\deleted{\unsure{I}} 1.4.2) that $\handwritten{\varphi}_i$ is a closed immersion, which was to be proved.}

\uline{Lemma 11.12.- Let $k$ be a field, $G$ a quasi-compact $k$-group. Put $A=\handwritten{\unsure{\mathscr{A}|\mathscr{O}}}(G)$. Given $x\in A$, there exists \deleted{a finite-dimensional $k$-vector subspace}\inserted{a $k$-subalgebra of finite type.} $V$ of $A$ such that $x\in V$, that $\handwritten{\Delta}(V)\subset V\handwritten{\otimes}_k V$ and $u(V)\subset V$, where $u$ denotes the involution of $A$ corresponding to inversion on $G$}

\inserted{By (11.3), one can replace $G$ by $G_{\mathrm{af}}$, so one can suppose $G$ affine.}

Let $(\handwritten{\varphi}_i)$ be a basis of $A$. Then $\handwritten{\Delta}(x)=\handwritten{\sum\varphi}_i\handwritten{\otimes}a_i$ and $\handwritten{\Delta}(\handwritten{\varphi}_i)=\handwritten{\sum\varphi}_j\handwritten{\otimes}b_{ji}$; axiom (HA1) of (\deleted{Exposition} I 4.2) shows that $\handwritten{\sum\varphi}_i\handwritten{\otimes\Delta}(a_i)=\handwritten{\sum\Delta}(\handwritten{\varphi}_i)\handwritten{\otimes}a_i$, whence $\handwritten{\Delta}(a_i)=\handwritten{\sum}b_{ij}\handwritten{\otimes}a_j$. Let then $V$ be the \deleted{vector subspace}\inserted{$k$-subalgebra} of $A$ generated by the $b_{ij}$ and the $u(b_{ij})$. It is clear that $V$ is a $k$-algebra of finite type. Axiom (HA1) also shows that $\handwritten{\sum\Delta}(\handwritten{\varphi}_j)\handwritten{\otimes}b_{ji}=\handwritten{\sum\varphi}_j\handwritten{\otimes\Delta}(b_{ji})$, and one deduces that $\handwritten{\Delta}(b_{ij})=\handwritten{\sum}b_{ik}\handwritten{\otimes}b_{kj}$ and that $\handwritten{\Delta}(u(b_{ij}))=\handwritten{\sum}u(b_{kj})\handwritten{\otimes}u(b_{ik})$. Since $\handwritten{\Delta}$ is an algebra homomorphism, one deduces that $\handwritten{\Delta}(V)\subset V\handwritten{\otimes}_k V$ and $u(V)\subset V$. Finally axiom (HA2) of (\deleted{Exposition} I 4.2) shows that $a_i=\handwritten{\sum\eta}(a_j)b_{ij}$ and that $x=\handwritten{\sum\eta}(\handwritten{\varphi}_i)a_i$, so that $x\in V$.

\uline{Proposition 11.13.- \deleted{Let $G$ be an affine $S$-group} Let $k$ be a field and $G$ an affine $k$-group with algebra $A$. . Then $G$ is the projective limit of an upward filtering system of affine $k$-groups of finite type, whose transition morphisms are faithfully flat}.

Since the set of $k$-subalgebras of finite type of $A$ stable under $\handwritten{\Delta}$ and $u$ is stable under generation, $A$ is the inductive limit of an upward filtering family $(B_i)_{i\in I}$ of $k$-subalgebras of finite type stable under $\handwritten{\Delta}$ and $u$. Then, each algebra $B_i$, equipped with the ho-\linebreak momorphism $B_i\handwritten{\longrightarrow}B_i\handwritten{\otimes}_k B_i$ deduced from $\handwritten{\Delta}$ and with the restriction of $u$ to $B_i$ is equipped with an associative augmented involutive hyperalgebra s\unsure{\deleted{rt}\inserted{tr}|\deleted{tr}\inserted{rt}}ucture, hence (\deleted{Exposition} I 4.2) is the \begin{tikzpicture}[x=1cm,y=1cm,baseline=(n.base)]\node (n) at (0,0) {$k$};\draw (-.12,.3)--(-.01,.38)--(.03,.4);\end{tikzpicture}-algebra of an affine $k$-group $G_i$, of finite type over $k$. Finally, since $A=\underset{\handwritten{\longrightarrow}}{\lim}V_{\begin{tikzpicture}[x=1cm,y=1cm,baseline=(n.base)]\node[inner sep=0pt] (n) at (0,0) {\scriptsize $i$};\draw (.01,-.04)--(.19,-.04)--(.15,-.09);\end{tikzpicture}}$,
% The top of one further character is visible at the clipped lower edge.
$\unsure{G}$
